\section{Estado del Arte}

\noindent El problema de localización de instalaciones es una línea clásica de estudio dentro de la investigación de operaciones, la optimización combinatoria y el diseño de redes logísticas. Su propósito general es determinar en qué ubicaciones abrir instalaciones y cómo asignar clientes a ellas, de modo que la demanda sea satisfecha al menor costo posible. Dentro de esta familia de problemas, el Capacitated Facility Location Problem (CFLP) incorpora capacidades máximas en las instalaciones, lo que permite representar sistemas donde las bodegas, plantas o centros de distribución no pueden atender una cantidad ilimitada de demanda \cite{klose2005facility}.\\

\noindent Históricamente, el CFLP ha sido abordado mediante modelos de programación matemática, especialmente formulaciones de programación lineal entera mixta. Estas representaciones consideran decisiones de apertura de instalaciones, asignación de clientes y uso de capacidad, permitiendo obtener soluciones óptimas en instancias pequeñas o medianas. Entre los enfoques exactos más relevantes se encuentran métodos de ramificación, planos de corte y descomposición de Benders. Sin embargo, la dificultad computacional aumenta significativamente cuando se incorporan restricciones adicionales, como capacidades ajustadas, asignación única o incompatibilidades entre clientes \cite{fischetti2016benders,klose2005facility}.\\

\noindent El Capacitated Facility Location Problem with Customer Incompatibilities (CFLP-CI) surge como una extensión del CFLP clásico al incorporar restricciones de incompatibilidad entre clientes. En esta variante, ciertos pares de clientes no pueden ser atendidos por una misma instalación, lo que permite modelar situaciones donde algunos productos, servicios o requerimientos no deben coexistir en una misma bodega por razones sanitarias, técnicas, operacionales o de seguridad. Esta característica modifica la estructura del problema, ya que la factibilidad de una solución no depende únicamente de satisfacer demanda y respetar capacidades, sino también de evitar conflictos entre clientes asignados a una misma instalación \cite{maia2023metaheuristic,ceschia2024multineighborhood}.\\

\noindent Desde el punto de vista del modelamiento, el CFLP-CI suele representarse mediante formulaciones de programación entera mixta que combinan variables de apertura de instalaciones, variables de flujo o asignación de demanda y variables binarias auxiliares para indicar si un cliente es atendido por una determinada bodega. Además, las incompatibilidades pueden representarse como un grafo de conflictos, donde los nodos corresponden a clientes y las aristas indican pares incompatibles. Esta representación es efectiva porque permite traducir las incompatibilidades en restricciones que impiden que dos clientes conectados por una arista sean servidos por la misma instalación \cite{maia2023metaheuristic,mansini2022kernel,ceschia2024multineighborhood}.\\

\noindent Una de las primeras aproximaciones específicas al CFLP-CI fue desarrollada por Maia et al., quienes propusieron y compararon distintas técnicas metaheurísticas para resolver instancias del problema. Entre los enfoques considerados se encuentran métodos evolutivos, búsqueda local, procedimientos greedy y estrategias híbridas basadas en minería de datos. Sus resultados muestran que las metaheurísticas son una alternativa adecuada para abordar instancias de mayor tamaño, donde los métodos exactos pueden requerir tiempos computacionales elevados o no alcanzar soluciones óptimas dentro de límites razonables \cite{maia2023metaheuristic}.\\

\noindent Posteriormente, Mansini y Zanotti abordaron problemas de localización con incompatibilidades desde una perspectiva matheurística, considerando variantes single-source y multi-source. Su propuesta se basa en un enfoque de Two-Phase Kernel Search, que combina programación matemática con una estrategia de reducción del espacio de búsqueda. Este tipo de método resulta relevante porque conserva parte de la estructura formal del modelo, pero evita explorar todas las variables posibles, concentrándose en subconjuntos prometedores para obtener soluciones de buena calidad en menor tiempo \cite{mansini2022kernel}.\\

\noindent La variante de mayor interés para este informe es el Multi-Source Capacitated Facility Location Problem with Customer Incompatibilities (MS-CFLP-CI), donde un cliente puede recibir demanda desde más de una bodega. Esta flexibilidad permite aprovechar mejor la capacidad disponible y reducir costos de operación, pero también aumenta la complejidad, porque las incompatibilidades deben verificarse en cada bodega donde un cliente recibe una cantidad positiva de productos. En esta línea, el trabajo de Ceschia y Schaerf representa uno de los avances más relevantes, al proponer un algoritmo de Simulated Annealing multi-vecindario diseñado específicamente para esta variante. En síntesis, la tendencia actual apunta al uso de metaheurísticas especializadas y matheurísticas, ya que los métodos exactos siguen siendo útiles para formular el problema, obtener cotas y resolver instancias pequeñas, mientras que los enfoques híbridos permiten abordar instancias grandes y altamente restringidas con mayor eficiencia computacional \cite{ceschia2024multineighborhood,mansini2022kernel,fischetti2016benders}.